% Generated by scripts/assemble.py; corrections belong in source/corrections.json.
\examhead{179}{西安交通大学 2024 年数学分析试卷(回忆版)}{169}
% p170-b01; source PDF page 170
\par\Needspace{5\baselineskip}
\blockanchor{p170-b01}
\majorquestion{一}填空题. 每题6分, 共60分.\par
% p170-b02; source PDF page 170
\blockanchor{p170-b02}
\question{1}求  $\lim_{x\to0}\frac{e^{x}\sin x-x(1+x)}{x^{3}}=$  \_\_\_\_.\par
% p170-b03; source PDF page 170
\blockanchor{p170-b03}
\question{2}求  $\lim_{n\to\infty}(\ln n - [\ln n]) =$  \_\_\_\_.\par
% p170-b04; source PDF page 170
\blockanchor{p170-b04}
\question{3}已知  $f(x) = x + \ln x$ ,  $g(x)$  是  $f(x)$  的反函数, 求  $g''(x) =$  \_\_\_\_.\par
% p170-b05; source PDF page 170
\blockanchor{p170-b05}
\question{4}求  $\sum_{n=1}^{\infty}\frac{\sin nx}{n!}=$  \_\_\_\_.\par
% p170-b06; source PDF page 170
\blockanchor{p170-b06}
\question{5}已知  $f(x)=|x|^{\alpha}$ ，求  $\frac{\partial^{2}f}{\partial x_{1}^{2}}+\frac{\partial^{2}f}{\partial x_{2}^{2}}+\cdots+\frac{\partial^{2}f}{\partial x_{n}^{2}}=$  \_\_\_\_.\par
% p170-b07; source PDF page 170
\blockanchor{p170-b07}
\question{6}已知 $f(x, y) = \begin{cases} \frac{x^3}{x^2 + y^2}, & x^2 + y^2 \neq 0; \\ 0, & x^2 + y^2 = 0. \end{cases}$ 求在 (0,0) 点沿 ($\cos \alpha, \sin \alpha$) 方向的方向导数 $=$\par
% p170-b08; source PDF page 170
\blockanchor{p170-b08}
\question{7}求不定积分  $\int \sin(\ln x) \, dx =$ \_\_\_\_ .\par
% p170-b09; source PDF page 170
\blockanchor{p170-b09}
\question{8}$x^{2} + y^{2} = 2az$ 和 $x^{2} + xy + y^{2} = a^{2}$ 的交线的最大值为\par
% p170-b10; source PDF page 170
\par\Needspace{6\baselineskip}
\blockanchor{p170-b10}
\question{9}已知 $\Omega$ 是由\par
\nopagebreak[4]
% p170-b11; source PDF page 170
\blockanchor{p170-b11}
\displayformula{(0, 0, 1), (0, 1, 1), (1, 1, 1), (0, 0, 2), (0, 2, 2), (2, 2, 2)}
% p170-b12; source PDF page 170
\blockanchor{p170-b12}
组成的 $\mathbb{R}^3$ 的一个棱台, 那么 $\iiint_{\Omega} \frac{1}{y^2 + z^2} \, \mathrm{d}x \, \mathrm{d}y \, \mathrm{d}z = \underline{\quad}$.\par
% p170-b13; source PDF page 170
\blockanchor{p170-b13}
\question{10}已知 $S: \frac{x^2}{a^2} + \frac{y^2}{b^2} + \frac{z^2}{c^2} = 1$, 则 $\iint_{S} x \, \mathrm{d}y \, \mathrm{d}z + y \, \mathrm{d}z \, \mathrm{d}x + z \, \mathrm{d}x \, \mathrm{d}y =$\par
% p170-b14; source PDF page 170
\par\Needspace{5\baselineskip}
\blockanchor{p170-b14}
\majorquestion{二}解答题. 每题 15 分, 共 90 分.\par
% p170-b15; source PDF page 170
\blockanchor{p170-b15}
\question{11}已知数列 $\{x_{n}\}$ 满足 $x_{mn} \leq x_{m} + x_{n}, x_{n} > 0$, 证明: $\lim_{n \to \infty} \sqrt[n]{x_n}$ 存在.\par
% p170-b16; source PDF page 170
\blockanchor{p170-b16}
\question{12}证明: $\lim_{h \to 0^{+}, k \to 0^{+}} \frac{f(x_0 + h) - f(x_0 - k)}{k + h}$ 存在的充要条件为 $f$ 在 $x_0$ 处可导.\par
% p170-b17; source PDF page 170
\blockanchor{p170-b17}
\question{13}证明梅尔滕定理: 设级数 $\sum_{n=1}^{\infty} a_n$ 绝对收敛, $\sum_{n=1}^{\infty} b_n$ 收敛, 且 $\sum_{n=1}^{\infty} a_n = A$, $\sum_{n=1}^{\infty} b_n = B$, 则 $\sum_{n=1}^{\infty} c_n = AB$,\par
% p170-b18; source PDF page 170
\blockanchor{p170-b18}
其中  $c_{n}=a_{1}b_{n}+a_{2}b_{n-1}+\cdots+a_{n}b_{1}=\sum_{k=1}^{n}a_{k}b_{n-k+1}.$\par
% p170-b19; source PDF page 170
\blockanchor{p170-b19}
\question{14}已知 $f: \mathbb{R}^n \to \mathbb{R}^m$ 是一个 $C^1$ 类映射, 且满足 $f'(x) = 0$, 证明: $f$ 为常值映射.\par
% p170-b20; source PDF page 170
\blockanchor{p170-b20}
\question{15}证明: 若 $f(x)$ 在 $[a, b]$ 上连续, 那么 $f$ 在 $[a, b]$ 上有界.\par
% p170-b21; source PDF page 170
\par\Needspace{6\baselineskip}
\blockanchor{p170-b21}
\question{16}设 $a, b, T$ 是三个正实数, $a \neq b$. 证明:\par
\nopagebreak[4]
% p170-b22; source PDF page 170
\blockanchor{p170-b22}
\displayformula{\frac {1}{T} \int_ {- T} ^ {T} \frac {\sin a t \cos b t}{t}   \mathrm{d} t = \left\{ \begin{array}{l l} 1 + O \left(\frac {1}{T (a - b)}\right), & a > b; \\ O \left(\frac {1}{T (a - b)}\right), & a <   b. \end{array} \right.}
